\documentclass[11pt]{article}
\usepackage[utf8]{inputenc}
\usepackage[T1]{fontenc}
\usepackage{lmodern}
\usepackage[margin=1in]{geometry}
\usepackage{csquotes}
\usepackage[style=authoryear,backend=biber,sorting=nyt,maxbibnames=99,giveninits=true,uniquename=init]{biblatex}
\usepackage[hidelinks]{hyperref}
\addbibresource{references.bib}
\begin{document}
\section*{Conclusions}

Read as a whole, the frontier of subsurface-drainage research is defined
less by an absence of ideas than by an absence of the evidence, tools, and
coordination needed to convert ideas into transferable guidance. Across
every domain surveyed here---monitoring, modelling, network mapping,
nutrient and water-table management, edge-of-field treatment, salinity and
reuse, unintended consequences, and adoption---the field's most secure
statements are about the limits of its own evidence. Headline results rest
disproportionately on short records, single sites, and narrow
agro-ecological windows, and their authors caution explicitly against
premature generalisation. The recurring signal is therefore a precondition
rather than a conclusion: progress in each domain is bounded by the length,
density, spatial replication, and consistency of the underlying
observations \parencite{Rogger2017,Kladivko2001}. This makes the observational foundation a frontier in its own
right, and one on which all others depend.

Four converging priorities emerge from that shared diagnosis. The first is
coordinated, long-term, multi-scale data. The single most widely voiced need
is for records long enough to separate management effects from climatic
noise and from slow, directional change in soils and hydrology, dense and
year-round enough to resolve the non-growing-season and event-scale windows
in which losses concentrate, and geographically broad enough to test whether
site-specific findings hold across other soils, climates, crops, and
management regimes. Layered onto record length are harmonisation and
comparability: standardised methods, complete site characterisation,
shared databases, and openly reported campaigns, without which accumulating
studies cannot support the meta-analyses and typologies that would
generalise them. Tracers, isotopes, and high-frequency source
apportionment recur as the means to resolve the pathways that bulk
flow-and-concentration records leave ambiguous, and experimental designs
capable of detecting the effect sought---paired basins, nested scales,
and defensible analysis of the unreplicated field trials that dominate the
literature---form the bridge between a lengthening record and transferable
conclusions.

The second priority is strengthened and validated model components. Because
field trials cannot span every combination of practice, soil, and climate,
process-based drainage and water-quality models are expected to carry much
of the predictive burden; yet their credibility is constrained less by their
formulation than by the narrowness of the conditions under which they have
been tested. Validation across a wider envelope, the replacement of
empirical coefficients with functions of measurable site attributes, and the
process-level measurements needed to constrain the least mature routines---
denitrification and the nitrogen cycle, phosphorus and its dominant
preferential-flow pathway, cold-season dynamics, and coupled seepage and
groundwater exchange---are near-universal priorities \parencite{Youssef2006,Radcliffe2015}. So too are explicit
attention to identifiability and uncertainty, so that models are reported
with propagated bounds rather than as point predictions, and the disciplined
integration of machine learning and coupled model structures, whose own
validation requirements the field is only beginning to confront. A parallel
and equally unresolved gap is spatial: the physical drainage network is
largely unmapped, so uncertainty about the network itself propagates into
every downstream prediction, and translating instrumented-plot knowledge to
the watersheds where water-quality goals are actually set remains a genuinely
open methodological problem requiring non-invasive detection, shared
databases, intensively monitored bridging catchments, and mechanistic
treatment of channel-network, connectivity, and depressional processes \parencite{Amado2018,Naz2009}.

The third priority is bundled, targeted, and climate-resilient management.
Across nutrients, water-table control, drainage design, and edge-of-field
treatment, the levers are better described than optimised. No single practice
is a silver bullet; the decisive gains lie in combining in-field and
edge-of-field practices, sequencing them, and placing them where they matter
most---on the hydrologically active areas, critical source areas, and
connected hotspots that generate a disproportionate share of loads. The
common thread is a shift from passive, uniform prescription toward active,
location-specific control: from one-way conveyance toward managed water
tables, automation, and forecast-based operation; from agronomic nutrient
rates toward site-specific crediting of legacy and mineralisable sources and
management targeted to the temporal windows of loss; from demonstration-scale
treatment structures toward standardised, transferable design across
climates and maturation stages \parencite{Randall2001,Strock2011,Christianson2021}. Underlying all of these is climate
adaptation: cropping and drainage systems, and the design criteria that
govern them, must function under both drought and flood, and must be
re-evaluated against long records rather than the short-term simulations on
which much current guidance rests.

The fourth priority is interdisciplinary work that links mechanism, tools,
and viable adoption, and it is inseparable from a growing recognition of
consequences that fall outside the nitrate- and yield-centred accounting
that has historically justified drainage. Removing a nutrient from one
pathway can shift it to another; treating drainage water can manufacture new
by-products; and altering the field water regime can set slow biogeochemical
and pedological change in motion. Greenhouse-gas coupling, pollution
swapping, emerging contaminants, trace-element legacies, and
drainage-induced soil evolution share a common feature: few studies have
quantified more than one consequence at a time, so the net balance among
competing objectives remains largely conjectural, and their full costs
cannot be read from a few years of outlet monitoring \parencite{Mitchell2023,Kadlec2012}. Resolving these
tradeoffs, and translating frontier science into practice, ultimately
depends on the human dimension---economics, incentives, policy,
institutions, and farmer engagement---and on integrative landscape-scale
frameworks and shared decision-support infrastructure that couple mechanistic
understanding with spatial targeting, cost-effectiveness, and stakeholder
knowledge \parencite{McLellan2015,Christianson2014}. Biophysical insight alone is necessary but not sufficient; without
economically and institutionally adoptable solutions, better science does not
become better outcomes.

These four strands are mutually reinforcing rather than independent.
Longer, harmonised data are what allow model components to be validated and
their uncertainty bounded; validated, coupled models are what allow
field-scale knowledge to be scaled to watersheds and management to be
targeted and made climate-resilient; and targeted, multi-objective
management is what makes the consequences worth quantifying and the
economics worth analysing. The unifying message of this review is therefore
that the next decade of progress in subsurface drainage will be secured less
by new equations or novel practices than by the sustained, well-instrumented,
openly reported, and geographically diverse field campaigns needed to test
the ideas already in hand, by the coordinated data infrastructure and
standards that make those campaigns cumulative, and by the interdisciplinary
integration that ties mechanism to tools and tools to adoption. The task
ahead is one of consolidation and coordination at scale---an agenda in which
data, models, management, and human-dimension advances must progress
together, because each is bounded by the others.

\printbibliography
\end{document}
